\chapter{Salud}

\marginal{el titular y lo que hay debajo}
La función Salud del Gobierno Federal más las entidades de control directo, en
bruto, pasa de \textbf{970,522.5 a 888,903.7 millones de pesos}: una caída real de
\textbf{12.2\%} y de 2.82 a \textbf{2.46\% del PIB}. Es el movimiento funcional
más grande del paquete.

\textbf{Y en su mayor parte no es un recorte al servicio.} Este capítulo lo
descompone antes de comentarlo, porque el titular es cierto de la función y falso
de la atención médica.

\section{Declaración de perímetro y de sustitución de fuente}

El perímetro canónico de esta casa para salud se ancla al cuadro ``Gasto en salud:
dependencias, entidades y ramos generales'' de la exposición de motivos del
proyecto de presupuesto. \textbf{Ese documento no existe para el ejercicio 2025}:
está caído en el servidor de la Secretaría, con las páginas índice enumerándolo y
los archivos devolviendo error.

El perímetro se reconstruye desde los analíticos por función, respetando los dos
bloques del canon. \textbf{La sustitución cambia el objeto y hay que decirlo}: los
analíticos son brutos y aquel cuadro era neto, de modo que las cifras de este
capítulo no son comparables con las que esta casa publicó en ejercicios anteriores
sin ajustar por el neteo.

\begin{table}[htbp]
\centering
\caption{Función Salud por ramo y entidad, en bruto (mdp)}
\label{tab:C6_1}
\small
\begin{tabular}{lSSSSS}
\toprule
{Concepto} & {2024} & {2025} & {Dif. nominal} & {Dif. real} & {Var. real \%} \\
\midrule
{TOTAL funcion Salud (GF + entidades, bruto)} & \num{970522.5} & \num{888903.7} & \num{-81618.8} & \num{-123351.3} & \num{-12.2} \\
\multicolumn{6}{l}{\itshape Gobierno Federal, por ramo} \\
{47 Entidades no Sectorizadas} & \num{128623.9} & \num{165352.2} & \num{36728.3} & \num{31197.5} & \num{23.3} \\
{33 Aportaciones Federales para Entidades Federativas y Municipios} & \num{135589.4} & \num{81220.5} & \num{-54369.0} & \num{-60199.3} & \num{-42.6} \\
{12 Salud} & \num{88916.7} & \num{59089.7} & \num{-29827.1} & \num{-33650.5} & \num{-36.3} \\
{19 Aportaciones a Seguridad Social} & \num{74226.5} & \num{9896.4} & \num{-64330.0} & \num{-67521.8} & \num{-87.2} \\
{07 Defensa Nacional} & \num{8009.4} & \num{7053.5} & \num{-955.8} & \num{-1300.2} & \num{-15.6} \\
{13 Marina} & \num{3659.0} & \num{3521.2} & \num{-137.8} & \num{-295.1} & \num{-7.7} \\
\multicolumn{6}{l}{\itshape Entidades de control directo} \\
{GYR Instituto Mexicano del Seguro Social} & \num{455245.4} & \num{481083.0} & \num{25837.6} & \num{6262.1} & \num{1.3} \\
{GYN Instituto de Seguridad y Servicios Sociales de los Trabajadores del Estado} & \num{76252.2} & \num{81687.2} & \num{5435.0} & \num{2156.1} & \num{2.7} \\
\bottomrule
\end{tabular}
\par\vspace{2pt}\footnotesize Fuente: elaboración propia del ITED con los analíticos del PEF aprobado 2024 y 2025. Las diferencias en millones de pesos se reportan dos veces, nominal y real; la real está en pesos de 2025 con el deflactor de 4.3\% que declara el CGPE.
\end{table}


\section{Los cuatro movimientos}

La caída de 81,618.8 millones de pesos nominales se explica por cuatro cosas de
naturaleza distinta.

\begin{itemize}
\item \textbf{Un adeudo que se liquidó y no se repite.} El programa de adeudos
con el IMSS y el ISSSTE del Ramo 19 pasa de 52,324.5 a 339.8 millones de pesos.
Fue una liquidación en 2024. Su desaparición no reduce la atención de nadie.
\item \textbf{Un programa que cambió de ramo.} IMSS-Bienestar sale del Ramo 19,
donde valía 21,623.7, y el organismo en Entidades no Sectorizadas sube de
128,623.9 a 165,352.2, un aumento real de 23.3\%.
\item \textbf{El fondo de aportaciones para los servicios de salud cae de verdad},
de 135,589.4 a 81,220.5 millones: 42.6\% real. Es la decisión de gasto más grande
del capítulo y la que merece el nombre de recorte.
\item \textbf{El Ramo 12 cae 36.3\% real}, de 88,916.7 a 59,089.7, en buena parte
por la misma reasignación institucional.
\end{itemize}

\marginal{$-12.2$ contra $-7.2$}
\textbf{Descontando solo el adeudo liquidado, la caída real pasa de 12.2 a 7.2\%.}
Es la cifra que un lector debería llevarse, con la advertencia de que sigue siendo
una caída considerable y que su origen está identificado: el fondo federalizado.

\section{Quién presta el servicio, sube}

Los dos institutos que atienden a la mayor parte de la población derechohabiente
\textbf{crecen} en términos reales: el IMSS 1.3\% (455,245.4 a 481,083.0) y el
ISSSTE 2.7\% (76,252.2 a 81,687.2). Entre los dos son el 63.3\% de la función en
2025.

La composición del gasto en salud se desplaza, por tanto, del componente
federalizado y de la Secretaría hacia los institutos y hacia IMSS-Bienestar.
\textbf{Quien pierde es la población sin seguridad social atendida por los estados
a través del fondo de aportaciones.}

\section{Comparaciones estructurales}

Las dos razones de la casa para salud dicen cosas distintas y conviene leerlas
juntas.

\emph{Gasto en salud del IMSS sobre cuotas obrero-patronales}: 0.85 en 2021, 0.851
en 2024 y \textbf{0.798 en 2025}. La razón fue notablemente estable durante cuatro
años y en 2025 cae cinco centésimas. Las cuotas crecen 12.7\% nominal y la salud
del instituto 5.7\%: \textbf{el instituto recauda más de lo que gasta en atender},
y la diferencia financia otras cosas, principalmente pensiones.

\emph{Pensiones sobre salud dentro del IMSS y el ISSSTE}: 2.11 en 2021, 2.282 en
2024 y \textbf{2.412 en 2025}. Por cada peso que los dos institutos gastan en
atender enfermos, gastan 2.41 en pagar pensiones, contra 2.11 hace cuatro años.
\textbf{Ésa es la tensión de fondo del sistema y ninguna pieza del paquete la
nombra.}

\section{El hueco que este capítulo no puede llenar}

La pregunta ``cuánto por afiliado'' no se responde. \textbf{No hay en el paquete
ni en la carpeta de este documento una fuente de población por afiliación}, y
dividir el gasto entre una población inventada sería peor que no dividirlo.

Se deja constancia de lo que haría falta: el padrón de derechohabientes del IMSS y
del ISSSTE y la población sin seguridad social estimada por el Consejo Nacional de
Población. Con eso, y no antes, el capítulo podría decir si el desplazamiento
hacia los institutos mejora o empeora el gasto por persona atendida.

\clearpage
